\section{Dependent types, with examples}\label{sec:01-basics:03-dependent-types:1}

Every function seen so far has a \emph{fixed} codomain. \texttt{double : Nat → Nat} returns a \texttt{Nat} regardless of the input \texttt{n}. The Lean type theory allows something more general, a type that itself depends on a \emph{value}, and a function whose return type changes depending on which argument it was given. This is a \textbf{dependent type}. It is the single feature that separates Lean (and other proof assistants) from an ordinary typed programming language, so it is worth building up slowly, from the smallest possible example, before naming the general pattern.

\subsection{The problem, in Python first}\label{sec:01-basics:03-dependent-types:2}

Before any Lean syntax, here is the actual problem dependent types solve, in a language with no type checker enforcing anything at all. Suppose a dot-product function is written over two lists:

\begin{lstlisting}[language=Python,style=python]
def dot(xs, ys):
    return sum(x * y for x, y in zip(xs, ys))

dot([17, -3, 42], [99, 8, 6])   # 1911 (*\textemdash{}*) correct
dot([17, -3, 42], [99, 8])      # 1659 (*\textemdash{}*) silently wrong: zip() truncates to the shorter list
\end{lstlisting}

This is worse than a crash. \texttt{dot([17, -\allowbreak{}3, 42], [99, 8])} does not raise anything. \texttt{zip} just quietly drops the third element of \texttt{xs} and hands back a number that \emph{looks} like a perfectly good answer. The bug (calling \texttt{dot} on two lists of different lengths, which is mathematically nonsensical) is never caught, not by Python, not by a test that happens not to cover this call site, not by anything short of a human noticing the number looks off.

The underlying mistake is that a Python \texttt{list} carries no information about its own length in any way the function signatures of Python can see or enforce. \texttt{xs : list} and \texttt{ys : list} say nothing about whether the two lists agree in size. That fact exists only as a comment at best, never checked by anything. What is actually wanted is a signature saying ``\texttt{dot} accepts two lists \emph{of the same length} \texttt{n}, for \emph{any} \texttt{n},'' a signature that mentions and constrains a value (\texttt{n}), not just a type (\texttt{list}). Ordinary type systems, including that of Python, even with type hints added, have no way to say that. This is exactly the gap dependent types close, and the rest of this section builds the machinery to say it, starting from the smallest possible example.

\subsection{First: a type family, one type per number}\label{sec:01-basics:03-dependent-types:3}

Here is the smallest genuine example of a dependent type, and it is already sitting in the core library of Lean. \href{https://leanprover-community.github.io/mathlib4_docs/Init/Prelude.html\#Fin}{\texttt{Fin n}} is the type of natural numbers \emph{strictly less than} \texttt{n}:

\begin{lstlisting}
#check Fin 3   -- Fin 3 : Type
#check Fin 5   -- Fin 5 : Type
\end{lstlisting}

\texttt{Fin} itself is not one type. It is a \textbf{recipe that produces a type once handed a number}. \texttt{Fin 3} and \texttt{Fin 5} are both genuine types, but they are \emph{different} types. \texttt{Fin 3} has exactly three inhabitants (the numbers \texttt{0}, \texttt{1}, \texttt{2}) and \texttt{Fin 5} has exactly five. Compare this to something already known not to be dependent, \texttt{List α}. \texttt{List Nat} and \texttt{List Bool} are different types too, but only because \texttt{Nat} and \texttt{Bool} are different \emph{types} fed in for \texttt{α}. \texttt{Fin} is a different kind of thing. \texttt{Fin 3} and \texttt{Fin 5} differ even though \texttt{3} and \texttt{5} are both perfectly ordinary terms of the exact same type, \texttt{Nat}. The type \texttt{Fin} produces depends on which \emph{value} it is given, not just which type. That value-dependence is the entire definition of ``dependent type,'' and \texttt{Fin} is the simplest possible example of one.

The construction of \texttt{Fin} can be inspected directly:

\begin{lstlisting}
#print Fin
-- structure Fin (n : Nat) : Type
-- fields:
--   Fin.val  : Nat
--   Fin.isLt : ↑self < n
\end{lstlisting}

So a term of \texttt{Fin n} is literally a pair, a \texttt{Nat} value, plus a \emph{proof} that the value is below \texttt{n}. The very statement of the proof (\texttt{↑self < n}, i.e.~the wrapped \texttt{val} compared against \texttt{n}) mentions \texttt{n}, the value supplied. Change \texttt{n} and the result is a genuinely different type, with a genuinely different proof obligation attached. This bundling of data, plus a proof whose statement depends on that data, is the second half of the dependent-types story (formalized later in this section as a \textbf{Σ-type}). \texttt{Fin} is a real, live example of it, not a made-up one.

\subsection{\texorpdfstring{Second: a function whose \emph{return type} changes with its argument}{Second: a function whose return type changes with its argument}}\label{sec:01-basics:03-dependent-types:4}

Now for the companion idea, a \textbf{dependent function}, one whose return type depends on the specific \emph{value} of its argument, not just the type of the argument. Define fixed-length vectors from scratch, the standard first example in any dependent-type-theory course:

\begin{lstlisting}
inductive Vec ((*$\alpha$*) : Type) : Nat (*$\rightarrow$*) Type where
  | nil  : Vec (*$\alpha$*) 0
  | cons : (*$\alpha$*) (*$\rightarrow$*) Vec (*$\alpha$*) n (*$\rightarrow$*) Vec (*$\alpha$*) (n + 1)
\end{lstlisting}

Read this exactly like \texttt{Nat}'s two-constructor definition from the previous section, with one new ingredient. \texttt{Vec α} is not a single type, it is a \textbf{family of types indexed by a \texttt{Nat}}. \texttt{Vec α 0}, \texttt{Vec α 1}, \texttt{Vec α 2}, and so on, are all different types, one per length, and the length is tracked \emph{in the type itself}, not just at runtime. \texttt{nil} builds the unique length-\texttt{0} vector, and \texttt{cons} takes an element and a length-\texttt{n} vector and produces a length-\texttt{(n+1)} vector. The \texttt{n} used on both sides of the arrow in \texttt{cons} is the \emph{same} \texttt{n}, so the constructor itself enforces ``one longer than whatever it started with.''

One detail that the careful treatment in Section 2 of \texttt{\{α : Type\}} does not account for is that \texttt{n} is never bound anywhere in the \texttt{cons} line above, and yet the declaration elaborates. That is the \textbf{\texttt{autoImplicit}} setting of Lean, on by default, which turns an unbound lowercase identifier in the type of a declaration into an implicit argument automatically, here inserting \texttt{\{n : Nat\}} for you. The same thing happens with \texttt{α} in the definitions below, which is why \texttt{\#check @Vec.replicate} prints a \texttt{\{α : Type\}} binder that nothing in the source wrote. Mathlib and most substantial projects set \texttt{autoImplicit := false}, precisely so that a typo cannot silently become a new implicit argument. In such a project, write the binders out:

\begin{lstlisting}
inductive Vec ((*$\alpha$*) : Type) : Nat (*$\rightarrow$*) Type where
  | nil  : Vec (*$\alpha$*) 0
  | cons {n : Nat} : (*$\alpha$*) (*$\rightarrow$*) Vec (*$\alpha$*) n (*$\rightarrow$*) Vec (*$\alpha$*) (n + 1)
\end{lstlisting}

That version binds \texttt{n} explicitly rather than relying on the setting, which is what the Mathlib-style projects in Chapter 13 expect.

Here is a function that \emph{builds} one of these, and its type is the dependent-function payoff:

\begin{lstlisting}
def Vec.replicate (a : (*$\alpha$*)) : (n : Nat) (*$\rightarrow$*) Vec (*$\alpha$*) n
  | 0     => Vec.nil
  | n + 1 => Vec.cons a (Vec.replicate a n)

#check @Vec.replicate
-- @Vec.replicate : {(*$\alpha$*) : Type} (*$\rightarrow$*) (*$\alpha$*) (*$\rightarrow$*) (n : Nat) (*$\rightarrow$*) Vec (*$\alpha$*) n

#eval (Vec.replicate (-42 : Int) 3 : Vec Int 3)
-- Vec.cons (-42) (Vec.cons (-42) (Vec.cons (-42) (Vec.nil)))
\end{lstlisting}

\textbf{A third style of writing a \texttt{def} body, and why the signature ends in a bare \texttt{→}.} Every \texttt{def} up to this point (\texttt{double}, \texttt{average}, \texttt{identity}) named every argument inside \texttt{(...)} and gave one term after \texttt{:=}. \texttt{Vec.replicate} above does something visibly different for its last argument: \texttt{(n : Nat)} never appears to the left of \texttt{:=}, the signature itself ends in a plain \texttt{→} the way the statement of a \texttt{theorem} does, and the body is a list of \texttt{| pattern => term} equations instead of one term. This style of writing a function body, naming the \emph{shape} an argument must have rather than a name for the argument itself, is called \textbf{pattern matching}, a term used again from here on wherever a definition or proof picks apart a value by its constructor rather than by a bound name. This is not a new arrow meaning on top of the two flagged in \hyperref[sec:01-basics:02-def-let-implicit:1]{Section 2}, and it is not a different kind of function. It is a third \emph{writing style} for the exact same thing, worth seeing side by side on a smaller example before trusting it in \texttt{Vec.replicate}:

\begin{lstlisting}
def vecLen {(*$\alpha$*) : Type} {n : Nat} (_ : Vec (*$\alpha$*) n) : Nat := n

def vecLen' {(*$\alpha$*) : Type} {n : Nat} : Vec (*$\alpha$*) n (*$\rightarrow$*) Nat
  | _ => n
\end{lstlisting}

\begin{itemize}
\tightlist
\item
  \texttt{vecLen} is the familiar style. The \texttt{Vec α n} argument is bound by name (as \texttt{\_\allowbreak{}}, since the body never actually uses it, only its index \texttt{n}) inside \texttt{(...)}, and the body is one term, \texttt{n}, exactly like the body of \texttt{double} was one term, \texttt{n * 2}.
\item
  \texttt{vecLen'} binds nothing after the colon. Its signature is \texttt{Vec α n → Nat}, a bare arrow, precisely as if it were the \emph{statement} of a theorem rather than the header of a function. The argument is instead supplied by the single equation \texttt{| \_\allowbreak{} => n} below, matching against whatever term of \texttt{Vec α n} is eventually passed in. Lean's \textbf{equation compiler} is what turns this \texttt{| pattern => term} block into an ordinary function, internally no different from writing \texttt{fun v => match v with | \_\allowbreak{} => n} inside a \texttt{:=} body. \texttt{vecLen} and \texttt{vecLen'} are, after elaboration, the same function under two different pieces of surface syntax, not two functions that happen to agree on every input.
\item
  The two styles are not interchangeable in every situation, only equally valid where they overlap. Equation-style earns its keep exactly when the whole purpose of a definition is to case on the \emph{shape} of an inductive argument, one equation per constructor, which is why \texttt{Vec.replicate} above (and \texttt{Vec.head}, \texttt{Vec.dot} shortly) are written that way. Each genuinely has one case for \texttt{0}/\texttt{nil} and one for \texttt{n + 1}/\texttt{cons}. The named-binder \texttt{:=} style stays the natural choice once the argument is used as an ordinary value rather than taken apart, exactly as \texttt{vecLen} above never inspects \emph{which} vector it received, only its already-known length \texttt{n}.
\end{itemize}

\begin{mathreading}

Equation-style \texttt{def} is the Lean transcription of ordinary definition by cases, \[
f(n) = \begin{cases} c_0 & n = 0 \\ c(k, f(k)) & n = k + 1, \end{cases}
\] the exact shape a mathematician already writes for a recursively defined sequence or function. Nothing new is being learned here beyond the Lean spelling of a pattern already familiar from ordinary practice, which is also exactly why it shows up unannounced the moment this book starts defining functions over an inductive type like \texttt{Vec}.

\end{mathreading}

\texttt{Vec.replicate} recurses exactly once per unit of length, so it is worth watching the recursion unwind one call at a time. Adding a \texttt{dbg\_\allowbreak{}trace} line to each branch (harmless, since it only prints, and changes nothing about what the function returns) makes every step visible:

\begin{lstlisting}
def Vec.replicate' (a : (*$\alpha$*)) : (n : Nat) (*$\rightarrow$*) Vec (*$\alpha$*) n
  | 0     => dbg_trace s!"replicate: n=0, base case, returning Vec.nil"; Vec.nil
  | n + 1 => dbg_trace s!"replicate: n={n+1}, prepending one copy of a, recursing with n={n}";
             Vec.cons a (Vec.replicate' a n)

#eval (Vec.replicate' (-42 : Int) 3 : Vec Int 3)
-- replicate: n=3, prepending one copy of a, recursing with n=2
-- replicate: n=2, prepending one copy of a, recursing with n=1
-- replicate: n=1, prepending one copy of a, recursing with n=0
-- replicate: n=0, base case, returning Vec.nil
-- Vec.cons (-42) (Vec.cons (-42) (Vec.cons (-42) (Vec.nil)))
\end{lstlisting}

Read the trace bottom-to-top against top-to-bottom. The four \texttt{dbg\_\allowbreak{}trace} lines print in the order each recursive call is \emph{entered} (\texttt{n=3}, then \texttt{n=2}, then \texttt{n=1}, then the base case \texttt{n=0}), and only once the base case returns does the final \texttt{\#eval} line print the fully-built \texttt{Vec}. This is the general shape every structurally-recursive function over \texttt{Vec}/\texttt{Nat} in this book has. One line prints per recursive call, in call order, with the line for the base case printed last before the result appears. The value of \texttt{a} itself cannot be printed generically inside the trace. \texttt{α} is an arbitrary type here, and nothing says a \texttt{Repr α} instance (the typeclass that lets a value be turned into displayable text) exists for whichever type \texttt{α} turns out to be at a given call site. This is why the message names the \emph{step} (which branch fired, what \texttt{n} is) rather than the data, the same reason \texttt{dbg\_\allowbreak{}trace} traces later in this book stick to naming steps and concrete, already-known values, never a still-generic argument.

This example deliberately stores \texttt{Int} values, not \texttt{Nat} values. With a \texttt{Vec Nat n} holding \texttt{Nat} elements, the length \texttt{n} and the stored numbers would both be \texttt{Nat}, and it becomes easy to lose track of which \texttt{Nat} is playing which role. \texttt{Int} keeps the elements visibly numeric, unlike, say, \texttt{Bool}, while still being unmistakably a \emph{different} type from \texttt{Nat}, the type of the length. A stored element such as \texttt{-\allowbreak{}42} can be large and negative, which a length can never be, so no reader could mistake it for the length \texttt{3}.

Real or complex numbers would make the distinction just as clear, but are not an option this early. \texttt{ℝ}/\texttt{ℂ} are Mathlib types, and this book stays Mathlib-free through Chapter 11. The reals in Lean in particular are \texttt{noncomputable} (built from Cauchy sequences with no decidable equality), so \texttt{\#eval} cannot evaluate one at all, not even in principle. \texttt{Int} is the closest numeric type that is both core Lean and actually computable. The next example, \texttt{Vec.dot}, uses \texttt{Vec Int n} for the same reason.

Look closely at the type \texttt{(n : Nat) → Vec α n}. The \texttt{n} that appears on the \emph{left} of the arrow (the argument) reappears inside the type on the \emph{right} of the arrow (the result). Feed \texttt{Vec.replicate a} the number \texttt{3}, and the result has type \texttt{Vec α 3}, specifically. Feed it \texttt{5}, and the very same function returns something of type \texttt{Vec α 5} instead. This is categorically different from \texttt{double : Nat → Nat}, where the output type (\texttt{Nat}) is fixed in advance and never reads the input value at all. Here, the \emph{type itself} changes based on which number was passed in. That is exactly what ``the codomain depends on the argument'' means, made as concrete as possible.

\subsection{Why bother: invariants become part of the type, not a side promise}\label{sec:01-basics:03-dependent-types:5}

The payoff is not just bookkeeping. Because the length lives in the type, Lean can rule out a whole class of mistakes \emph{before running anything at all}. Define a function that reads the first element of a vector, which only makes sense for a \emph{non-empty} vector:

\begin{lstlisting}
def Vec.head : Vec (*$\alpha$*) (n + 1) (*$\rightarrow$*) (*$\alpha$*)
  | Vec.cons a _ => a
\end{lstlisting}

The argument type \texttt{Vec α (n + 1)} says, in the type itself, ``this only accepts vectors of length \emph{at least one}.'' There is no separate runtime check for emptiness anywhere in this definition, because none is needed.

Notice also that the equation-style body above has only one case, \texttt{Vec.cons a \_\allowbreak{} => a}, with no \texttt{Vec.nil} arm. In other words, \texttt{Vec.nil} is not a case this match forgot or silently skips. It is a case Lean will not even let be \emph{written} here. \texttt{Vec.nil} has type \texttt{Vec α 0}, but the argument being matched has type \texttt{Vec α (n + 1)}, so a \texttt{Vec.nil} arm would require \texttt{0} to unify with \texttt{n + 1} for some \texttt{n}, and no such \texttt{n} exists, since \texttt{Nat.zero} and \texttt{Nat.succ n} are different constructors of the same type and Lean already knows two different constructors can never produce the same value. The equation compiler checks exactly this while elaborating the match, sees the \texttt{Vec.nil} case is impossible at this type, and closes it off automatically. What looks like an \emph{incomplete} match, only one pattern for a type with two constructors, is actually a \emph{complete} one, once the index \texttt{n + 1} is taken into account.

Calling it on an empty vector is rejected before the expression ever runs:

\begin{lstlisting}
#check Vec.head Vec.nil
\end{lstlisting}

\begin{lstlisting}
error: Application type mismatch: The argument
  Vec.nil
has type
  Vec ?m 0
but is expected to have type
  Vec ?m (?n + 1)
in the application
  Vec.nil.head
\end{lstlisting}

The pretty-printer in Lean renders the failed application in field-notation form, \texttt{Vec.nil.head}, rather than echoing back \texttt{Vec.head Vec.nil} verbatim. This is cosmetic, but worth not being surprised by when reproducing this error directly.

Nothing about ``index out of range'' happens at runtime, because the bad call is not a well-typed term in the first place. This is the same ``ruled out before running'' guarantee from \hyperref[sec:01-basics:01-everything-has-a-type:1]{Chapter 1, Section 1}, now enforced by an invariant (non-emptiness) that an ordinary, non-dependent type could not have expressed at all. \texttt{List α} has no way to say ``and this one is non-empty'' as part of its type. \texttt{Vec α (n+1)} says exactly that, for free, using only the machinery already on the table.

\textbf{Return to the Python example from the start of this section.} Here is \texttt{dot}, rewritten for \texttt{Vec} instead of the \texttt{list} type of Python, with both arguments required to share the \emph{same} length \texttt{n}:

\begin{lstlisting}
def Vec.dot : Vec Int n (*$\rightarrow$*) Vec Int n (*$\rightarrow$*) Int
  | Vec.nil, Vec.nil => 0
  | Vec.cons x xs, Vec.cons y ys => x * y + Vec.dot xs ys
\end{lstlisting}

The signature \texttt{Vec Int n → Vec Int n → Int} uses the \emph{same} \texttt{n}, a \texttt{Nat}, the length, for both arguments. That is not a naming coincidence, it is the whole point. Elements are \texttt{Int}, not \texttt{Nat}, on purpose. This keeps the length (\texttt{n}, a \texttt{Nat}) and the stored numbers (\texttt{Int}) as visibly different types throughout, with no \texttt{Nat}/\texttt{Nat} overlap left to lose track of. The vectors below are also named \texttt{vecA}/\texttt{vecB}, not after their own lengths, to avoid yet another coincidence layered on top. Try to reproduce the silent bug from Python:

\begin{lstlisting}
def vecA : Vec Int 3 := Vec.cons 17 (Vec.cons (-3) (Vec.cons 42 Vec.nil))
def vecB : Vec Int 2 := Vec.cons 99 (Vec.cons 8 Vec.nil)

#check Vec.dot vecA vecB
\end{lstlisting}

\texttt{Vec.dot} recurses on \emph{both} arguments at once, consuming one element from each per call, and accumulates a running total on the way back up. Before seeing why the mismatched-length call above fails, watch a same-length call succeed, with a \texttt{dbg\_\allowbreak{}trace} line added to each branch. Unlike the trace for \texttt{Vec.replicate}, this one can also print the actual numbers, since \texttt{Int} (unlike a fully generic \texttt{α}) does have a \texttt{Repr} instance:

\begin{lstlisting}
def Vec.dot' : Vec Int n (*$\rightarrow$*) Vec Int n (*$\rightarrow$*) Int
  | Vec.nil, Vec.nil => dbg_trace s!"dot: both nil, base case, returning 0"; 0
  | Vec.cons x xs, Vec.cons y ys =>
      dbg_trace s!"dot: heads x={x}, y={y}, recursing on the rest";
      x * y + Vec.dot' xs ys

def vecC : Vec Int 3 := Vec.cons 2 (Vec.cons 5 (Vec.cons 1 Vec.nil))

#eval Vec.dot' vecA vecC
-- dot: heads x=17, y=2, recursing on the rest
-- dot: heads x=-3, y=5, recursing on the rest
-- dot: heads x=42, y=1, recursing on the rest
-- dot: both nil, base case, returning 0
-- 61
\end{lstlisting}

Each traced line names the pair of heads being multiplied \emph{before} the recursive call is made, so the four lines above are exactly \(17 \times 2\), \(-3 \times 5\), \(42 \times 1\), and the \(0\) of the base case, in that order. These are the same four terms the recursion in the definition adds together, made visible one at a time instead of only appearing as the final sum \(34 + (-15) + 42 + 0 = 61\). \texttt{vecA} and \texttt{vecB} have different lengths on purpose, to set up the type-mismatch check next. \texttt{vecC} above is a third vector, the same length as \texttt{vecA}, used only to give \texttt{Vec.dot'} a legal pair to run on.

\begin{lstlisting}
error: Application type mismatch: The argument
  vecB
has type
  Vec Int 2
but is expected to have type
  Vec Int 3
in the application
  vecA.dot vecB
\end{lstlisting}

the version of Python, \texttt{dot([17,-\allowbreak{}3,42], [99,8])}, silently returned \texttt{1659}, a wrong answer with no error at all. The Lean version does not even compile. The length-mismatch bug is not caught by a clever runtime check \emph{added} to \texttt{Vec.dot}. There is no such check anywhere in its three-line definition. It is caught because ``both arguments have the same length'' was stated once, in the type, and Lean enforces every type it is given, automatically, for every call site, without exception.

The actual, built-in \texttt{Vector α n} in Lean is distinct from the toy \texttt{Vec} built here. It is defined differently under the hood, as an \texttt{Array α} paired with a proof that its size equals \texttt{n}, for performance reasons, the same way the \emph{presentation} of \texttt{Nat} above as \texttt{zero}/\texttt{succ}, due to Peano, does not reflect how Lean actually stores numbers at runtime (as fast arbitrary-precision integers). The \texttt{Vec} built in this section is the traditional textbook definition, simpler to reason about and the one every type-theory reference uses first, while the real \texttt{Vector} in Lean is engineered for speed. Both are dependent types in exactly the sense described here.

\subsection{The general pattern: Π-types}\label{sec:01-basics:03-dependent-types:6}

Both examples above are instances of one idea. A \textbf{dependent function type}, written with \(\Pi\) (``Pi,'' for ``dependent product''), generalizes the ordinary function-space \(A \to B\):

\[
\prod_{x : A} B(x)
\]

Here \(B\) is not itself a type. \(B\) is a \textbf{family of types indexed by \(A\)}, formally a function \(B : A \to \mathrm{Type}\) (or into \texttt{Prop}, the type of propositions, a distinct universe of its own, formally named \texttt{Sort 0} in \hyperref[sec:01-basics:05-pi-sigma-and-coc:1]{Chapter 1, Section 5}, as below). For each \(x : A\), \(B(x)\) is the specific type that family produces at \(x\), and different values of \(x\) may give genuinely different types. Read the whole expression as ``a function that, given any \(x : A\), returns a term of type \(B(x)\),'' a type allowed to mention \(x\), because \(B\) itself is allowed to vary with \(x\). When \(B(x)\) does not actually depend on \(x\) (i.e.~\(B\) is a \emph{constant} family), this collapses exactly to the ordinary function type \(A \to B\). Π-types \textbf{strictly generalize} function types. They do not replace them. The type of \texttt{Vec.replicate} above literally \emph{is} \(\prod_{n : \mathtt{Nat}} \mathrm{Vec}\,\alpha\,n\), with the surface syntax of Lean \texttt{(n : Nat) → Vec α n} spelling out the same thing without needing the \(\Pi\) symbol.

This is also, not by coincidence, exactly what \texttt{∀} means. \texttt{∀ n : Nat, n ≥ 0} is a Π-type where \(B(n)\) happens to be a \emph{proposition} (\texttt{n ≥ 0 : Prop}) rather than a data type like \texttt{Vec α n}. It reads as ``for every \texttt{n}, produce a proof of the \texttt{n}-specific statement \texttt{n ≥ 0}.'' Chapter 3 introduces \texttt{∀} and propositional logic properly. Once it does, every \texttt{∀} written from that point on is already a dependent function in exactly this sense, whether or not this vocabulary is available yet when it is first met. Propositions are just the special case where the family \(B\) happens to land in \texttt{Prop} instead of \texttt{Type}.

\subsection{Looking ahead}\label{sec:01-basics:03-dependent-types:7}

Chapter 11 builds a genuinely more elaborate dependent type, \texttt{Path Q : V → V → Type}, a family of types indexed by a \emph{pair} of vertices in a graph rather than by a single \texttt{Nat}. It is ``the type of paths from \texttt{u} to \texttt{w},'' which differs for each choice of endpoints exactly as \texttt{Vec α n} differs for each length. Its \texttt{cons}-like constructor is a dependent function for the same reason \texttt{Vec.replicate} is one here. Composing two paths is only accepted by the type-checker when their endpoints actually match, an invariant baked into the type rather than checked separately. Nothing new is needed to understand it once the \texttt{Fin}/\texttt{Vec} examples in this section make sense. It is the identical idea, with a richer index.

\begin{quote}
\textbf{Mathematical reading (optional).} For readers who already think categorically, an indexed family \texttt{B : A → Type} is exactly a functor out of the discrete category on \texttt{A}, or, thinking of \texttt{A × A}-indexed families as in the \texttt{Path} example above, an assignment of a \(\mathrm{Hom}\)-set to every pair of objects in a category. A Π-type over such a family is a \textbf{dependent product}. A term of \(\sum_{x:A} B(x)\) (Σ-type, next covered formally in \hyperref[sec:01-basics:05-pi-sigma-and-coc:1]{Chapter 1, Section 5}) is a \textbf{dependent sum}. Both are literal categorical limits/colimits in the appropriate indexed sense, not merely named after them by analogy.
\end{quote}

\begin{quote}
Read more. \hyperref[sec:01-basics:05-pi-sigma-and-coc:1]{Chapter 1, Section 5} gives Π-types (and Σ-types) their formal typing rules, with more worked examples, rather than only the walkthrough given here.
\end{quote}

\subsection{Sources, quoted}\label{sec:01-basics:03-dependent-types:8}

Formal definitions and citations for this section, gathered here for reference (full entries in the \hyperref[sec:root:bibliography:1]{Bibliography}):

\begin{itemize}
\tightlist
\item
  \textbf{Dependent type.} ``The short explanation is that types can depend on parameters'' (\cite{TPIL4}, §2.8 ``What makes dependent type theory dependent?''). Picture it like this. A cake sized to the guest count. ``The cake for 12'' and ``the cake for 3'' aren't the same cake made smaller, they're genuinely different objects, one per number of guests. \texttt{Fin n} is that idea turned into a type. \texttt{Fin 12} and \texttt{Fin 3} are different types, one per \texttt{n}, the same way cake sizes differ per guest count.
\item
  \textbf{Dependent function type (\(\Pi\)-type).} ``The type \texttt{(a : α) → β a} denotes the type of functions \texttt{f} with the property that, for each \texttt{a : α}, \texttt{f a} is an element of \texttt{β a}'' (\cite{TPIL4}, §2.8). Picture it like this. A vending machine whose dispensing slot changes shape depending which button is pressed. A soda comes out one opening, a bag of chips another. Written here as \(\prod_{x:A} B(x)\) (standard mathematical notation for the same construct; TPIL4 itself uses ``dependent function type''/``dependent arrow type,'' not the \(\Pi\) symbol), an ordinary function type \(A \to B\) is just the boring special case where every slot happens to be the same shape.
\item
  The Lean 4 source / {[}Mathlib4 API documentation{]}\cite{Mathlib4Docs} for \texttt{Fin} and \texttt{Vector}, confirmed directly in this section via \texttt{\#print Fin} against the actual toolchain pinned in \texttt{lean\_\allowbreak{}project/\allowbreak{}lean-\allowbreak{}toolchain} in this repository.
\item
  Thompson (\cite{Thompson1991}), §4.6 ``Quantifiers,'' §6.3 ``Dependent types and quantifiers'' develop the same dependent-product/dependent-sum content, verified verbatim against the source. The primary notation used by Thompson is \(\forall\)/\(\exists\), not Π/Σ. He calls the Σ-type-equivalent an ``(infinitary) sum type'' or ``dependent sum type,'' and the literal term ``Sigma-type'' never appears in his main text, only once, in a bibliography entry citing a paper by a different author. Explicit Π-notation does appear later, in the meta-theory chapters of that book (§8.3, §9.1.5), applied to dependent function spaces in a typed λ-calculus meta-language.
\item
  Chlipala (\cite{Chlipala2013}), §8.1 ``Length-Indexed Lists'' and §9.1 ``More Length-Indexed Lists''. The length-indexed-vector idea in this book (\texttt{ilist : nat → Set}) is built and revisited there, not in Ch. 2--3 as an earlier draft of this section stated, verified verbatim (\texttt{Inductive ilist : nat → Set := | Nil : ilist O | Cons : ∀ n, A → ilist n → ilist (S n)}). This is a useful second angle on the identical concept, in Coq rather than Lean.
\end{itemize}
